\section{Proofs for Section~\ref{sec:regimes}}\label{app:proofs-regimes}

Notation is as in Section~\ref{sec:regimes}. Throughout, $\mathcal A_r(T)=(1-e^{-rT})/r$, and for a cash-flow measure $\mathrm{CF}$ we write $\int e^{-rt}d\,\mathrm{CF}(t)$ for its present value. The script \texttt{code/verify\_part2\_regimes.py} checks each step marked below symbolically or on random draws.

\subsection*{Proof of Proposition~\ref{prop:reg-value}}

(a) Split the expected cash flows of Assumption~\ref{ass:tasks-transition} at $T_{\mathsf o}$. Because $T_{\rm par}\le T_{\mathsf o}$, every one-time cost, the parallel run and the stranded and rollback terms fall before $T_{\mathsf o}$. The proof of Proposition~\ref{prop:tasks-transition}(a) evaluates the expectation branch by branch (success, detected failure, undetected failure) and uses $T_h$ only as the date at which the savings stream ends. The flows on $[0,T_{\mathsf o})$ therefore have present value \eqref{eq:tasks-npv} at $r=r_{\mathsf o}$ with $T_h$ replaced by $\min\{T_h,T_{\mathsf o}\}$. At $T_{\mathsf o}$, on the success branch the run-rate cash saving is $bW$ and the run-rate EBITDA saving is $bW^{\rm E}$, because EBITDA does not charge the replacement capital spending $\delta_KI$ that $W$ nets out; both have been in place for at least $T_{\rm sea}$. On either failure branch the legacy mode runs at $T_{\mathsf o}$ and the run-rate saving is zero. Hence $\mathbb E\,\mathrm{TV}=h_xM_{\rm exit}\varpi_p\bar b_pW^{\rm E}$ under the EBITDA rule. Under the market rule, Lemma~\ref{lem:valuation2} gives $V=S_T[m\mathcal A_T(\eta)+(\theta m+(1-\theta)\bar m)\eta^T\Gamma]$, which is linear in $m$ with $\partial V/\partial m=S_T\mathfrak A_T$, so a permanent margin change $bW/S_T$ raises the price by exactly $\mathfrak A_TbW$, and $\mathbb E\,\mathrm{TV}=\mathfrak A_T\varpi_p\bar b_pW$. The last term is the definition \eqref{eq:reg-value} on the planned path. (Checks A4, G5.)

(b) For a going concern, $V=-C+W\int_t^\infty e^{-ru}du=-C+We^{-rt}/r$. Under an exit rule, $V=-C+W\int_t^Te^{-ru}du+e^{-rT}\widetilde MW\mathbf 1\{t\le T-T_{\rm sea}\}$, which is \eqref{eq:reg-K}. Let $c:=e^{-rT}(1-r\widetilde M)$, so that $\mathfrak K^{\rm x}(t)=(e^{-rt}-c)/r$ for $t\le T-T_{\rm sea}$ and $\mathfrak K^{\rm x}(0)=(1-c)/r>0$, because $c\le e^{-rT}<1$. Then
\[
\frac{\mathfrak K^{\rm x}(t)}{\mathfrak K^{\rm x}(0)}-e^{-rt}=\frac{e^{-rt}-c-(1-c)e^{-rt}}{1-c}=-c\,\frac{1-e^{-rt}}{1-c}.
\]
For $t>0$ the fraction is positive, so the difference has the sign of $-c$, which is positive if and only if $r\widetilde M>1$. For $t>T-T_{\rm sea}$ the terminal credit is lost. (Checks A1 to A3.)

(c) For a going concern, $\Delta(r):=V_\ell-V_q=-(C_\ell-C_q)+[W_\ell e^{-rt_\ell}-W_qe^{-rt_q}]/r$. Multiply by $re^{rt_q}>0$:
\[
h(r):=re^{rt_q}\Delta(r)=-re^{rt_q}(C_\ell-C_q)+W_\ell e^{-r(t_\ell-t_q)}-W_q .
\]
The first term is nonincreasing in $r$ because $C_\ell\ge C_q$, and the second is strictly decreasing because $W_\ell>0$ and $t_\ell>t_q$. As $r\downarrow0$, $h\to W_\ell-W_q>0$; as $r\to\infty$, $h\to-W_q<0$ if $C_\ell=C_q$ and $h\to-\infty$ otherwise. By continuity and strict monotonicity, $h$ has a unique root $r^\ddagger$, with $V_\ell>V_q$ below it and $V_\ell<V_q$ above it. Under an exit rule with $t_\ell>T-T_{\rm sea}$, part (b) removes the slow design's terminal credit. (Checks A5, A8.)

(d) At constant volume $v$, the switch spends $v\Delta K_j$ at date 0, saves $v(\Delta c_j-\delta_K\Delta K_j)$ a year in cash (replacement capital is a cash cost) and $v\Delta c_j$ a year in EBITDA. With the EBITDA rule its value is
\[
-v\Delta K_j+v(\Delta c_j-\delta_K\Delta K_j)\mathcal A_r(T)+e^{-rT}h_xM_{\rm exit}v\Delta c_j,
\]
which is positive if and only if $\Delta c_j[\mathcal A_r(T)+e^{-rT}h_xM_{\rm exit}]>\Delta K_j[1+\delta_K\mathcal A_r(T)]$, that is, \eqref{eq:reg-zeta}. For a public going concern without myopia, Proposition~\ref{prop:tasks-capital}(a) at $r_w$ gives the value $v[\Delta c_j-(r_w+\delta_K)\Delta K_j]/r_w$. A ratio between $\zeta^{\rm pe}$ and $r_w+\delta_K$ satisfies the first condition and fails the second. (Checks A9, A10.)

Finally, for the family saving in Assumption~\ref{ass:reg-regime}(v): by Definition~\ref{def:ext-effcost}, the owner values an own hour at $\hat w_j=w_j-\omega_h/(1-\tau)=(1-b_{\rm fam})w_j$. Redeployed staff keep their jobs, so only the displaced wage bill that is not redeployed, $(1-u_d)C^{\rm disp}$, is valued at $1-b_{\rm fam}$ rather than one, which lowers the saving by $b_{\rm fam}(1-u_d)C^{\rm disp}$. \qed

\subsection*{Proof of Proposition~\ref{prop:reg-weight}}

(a) Proposition~\ref{prop:dynamic} uses the discount rate only as a parameter with $r>G$, so it holds at $r=r_{\mathsf o}$. For the small-share weight, let $F(v;r):=(r-G)v-(m_0-kG)-C_\alpha[\lambda(v-k)_+]^{1/(1-\alpha)}$. Then $F(0;r)=-(m_0-kG)<0$ and, by \eqref{eq:gordon}, $v_0(r)$ is the smallest positive root, assumed simple. Since $F<0$ on $(0,v_0)$, $F_v(v_0;r)\ge0$, and simplicity gives $F_v>0$. $F$ is continuously differentiable near $(v_0,r)$ (if $v_0\le k$, $F$ is affine in $v$ there with slope $r-G>0$). Because $F$ is bounded away from zero below on $(0,v_0-\epsilon]$, small changes in $r$ keep $v_0$ the smallest root, and the implicit function theorem gives $dv_0/dr=-F_r/F_v=-v_0/F_v<0$. (Check B2.)

(b) With $m=m_0-\sigma-kg$ and $m_{op}=m_0-\sigma$,
\[
m+Xg+\mathfrak w m_{op}=(1+\mathfrak w)(m_0-\sigma)+(X-k)g=(1+\mathfrak w)\Bigl[m_0-\sigma+\frac{X-k}{1+\mathfrak w}g\Bigr]=(1+\mathfrak w)\,R_{X^{\rm my}}(g,m),
\]
because $R_{X^{\rm my}}(g,m)=X^{\rm my}g+m_0-\sigma-kg$. A positive multiple has the same maximizer, so Proposition~\ref{prop:score} applies with $y=X-k$ replaced by $y/(1+\mathfrak w)$. By \eqref{eq:Dstar} and \eqref{eq:sigstar}, $\Delta^*\propto y^{\alpha/(1-\alpha)}$ and $\sigma^*\propto y^{1/(1-\alpha)}$, which gives the factors $\rho$ and $(1+\mathfrak w)^{-1/(1-\alpha)}$. For the loss, $R_X(g,m)=R_\infty+y\Delta-\sigma$ on the frontier. At $\Delta'=\rho\Delta^*$ the spend is $\sigma'=(\Delta'/\ell)^{1/\alpha}=\rho^{1/\alpha}\sigma^*=\alpha\rho^{1/\alpha}y\Delta^*$, so $R_X-R_\infty=y\Delta^*(\rho-\alpha\rho^{1/\alpha})$, against $\Pi=(1-\alpha)y\Delta^*$ at the optimum. The difference is $\Pi[1-(\rho-\alpha\rho^{1/\alpha})/(1-\alpha)]$. (Checks B3 to B5.)

(c) Formally, as in Proposition~\ref{prop:dynamic}(c), the value $V(S,M,t)$ with terminal condition $V(\cdot,\cdot,T)=\mathrm{TV}$ satisfies $rV=\max_\sigma\{Sm+V_SSg+V_MMG+V_t\}$ with $g=G+\ell\sigma^\alpha$. The part that depends on $\sigma$ is $S[m_0-\sigma-kg+V_Sg]=S\,R_X(g,m)$ with $X=V_S$, and the argument does not use the horizon. A marginal dollar of revenue added at $t$ is $e^{G(u-t)}$ at $u\ge t$, earns $\hat me^{G(u-t)}$ until $T$ and is worth $X_Te^{G(T-t)}$ at $T$, so
\[
\frac{\partial V}{\partial S}(t)=\int_t^Te^{-(r-G)(u-t)}\hat m\,du+e^{-(r-G)(T-t)}X_T,
\]
which is \eqref{eq:reg-Xexit} with $\tau=T-t$. Under the EBITDA rule a dollar of revenue at $T$ adds $\hat m^{\rm E}$ of seasoned EBITDA, capitalized at $h_xM_{\rm exit}$; under the market rule $V_T$ is proportional to $S_T$ at given margin and growth (Lemma~\ref{lem:valuation2}), so $X_T=(V/S)_T$. Finally $dX^{\rm x}/d\tau=(r-G)e^{-(r-G)\tau}(X^{\rm G}-X_T)$ has a constant sign. (Checks B6, B7.)

(d) By \eqref{eq:Dstar}, $\Delta^*=\ell^{1/(1-\alpha)}[\alpha(X-k)]^{\alpha/(1-\alpha)}$ is increasing in $X>k$, as are $\sigma^*=\alpha(X-k)\Delta^*$ and $g^*$, so $m^*(g^*)=m_0-kg^*-\sigma^*$ is decreasing; $R^*_X$ is \eqref{eq:Rstar}. (Checks B8, B9.)

(e) The exact effect of $\delta_\sigma$ alone is $\Pi[(1-\delta_\sigma)^{-\alpha/(1-\alpha)}-1]$ by Proposition~\ref{prop:tasks-aggregate}(b), and $R^*_X$ is affine in $m_0$ with slope one (Proposition~\ref{prop:envelope}), so a cut $da$ raises it by exactly $da$. Since $\Pi(X)=C_\alpha[\ell(X-k)]^{1/(1-\alpha)}$ is strictly increasing for $X>k$, $\Pi(X)[(1-\delta_\sigma)^{-\alpha/(1-\alpha)}-1]>da$ if and only if $X-k>\ell^{-1}\{da/[C_\alpha((1-\delta_\sigma)^{-\alpha/(1-\alpha)}-1)]\}^{1-\alpha}$, which is \eqref{eq:reg-Xdag}. (Checks B10, B11.) \qed

\subsection*{Proof of Proposition~\ref{prop:reg-speed}}

(a) $\partial a/\partial r=[1-e^{-rT}-rTe^{-rT}]/(1-e^{-rT})^2=e^{-rT}(e^{rT}-1-rT)/(1-e^{-rT})^2>0$, since $e^x>1+x$ for $x>0$. If the alternatives carry no capital, $\underline c_j$ does not depend on $r$, while $c^{A,\mathrm{var}}_j=\pi n_j+c^Q_j+(r+\delta_K)K_j(A)$ and the build charge $F^A_ja(r,T_A)$ rise with $r$. Hence $v^*_j=F^A_ja/(\underline c_j-c^{A,\mathrm{var}}_j)$ rises where it is finite, the make-or-buy volume of Proposition~\ref{prop:tasks-window}(c) rises in the same way, and $\pi^*_j=\min_x[c^0_j(x)-c^0_j(A)]/[n_j-n_j(x)]$ falls because $c^0_j(A)$ rises. The capital test is Proposition~\ref{prop:tasks-capital}(a) at $r_{\mathsf o}$. In Proposition~\ref{prop:ext-timing}(b), the net flow $N(T)=v(T)s(T)-(r+\delta_F)F_A(T)$ is strictly decreasing in $r$ at each $T$, so if $N\ge0$ at some date for the rate $r'>r$, it is positive there for $r$, and the first such date at $r'$ is not earlier than at $r$.

For the internal rate of return, let the expected cash-flow measure be nonpositive on $[0,t^*]$ and nonnegative on $(t^*,\infty)$, and nonzero on the second set. Let $\Psi(r):=e^{rt^*}\mathrm{NPV}(r)=\int e^{-r(t-t^*)}d\,\mathrm{CF}(t)$. For $r<r'$, on $[0,t^*]$ we have $t-t^*\le0$, so $e^{-r(t-t^*)}\le e^{-r'(t-t^*)}$, and multiplying by $d\,\mathrm{CF}\le0$ reverses the inequality; on $(t^*,\infty)$, $e^{-r(t-t^*)}>e^{-r'(t-t^*)}$ and $d\,\mathrm{CF}\ge0$. Hence $\Psi(r)>\Psi(r')$: $\mathrm{NPV}$ changes sign at most once, from positive to negative, at the internal rate of return, and it is positive exactly below it. For the transition of Assumption~\ref{ass:tasks-transition}, the lump at date 0 and the expected stranded and rollback costs at $T_{\rm par}$ are negative, the flow on $[0,T_{\rm par})$ is $\varpi_p\bar b_pW\varrho(t)-C^{\rm par}\le\varpi_p\bar b_pW-C^{\rm par}\le0$, and the flow after $T_{\rm par}$ is $\varpi_p\bar b_pW\ge0$, so $t^*=T_{\rm par}$. Truncating the savings at any $T_h$ keeps the single sign change. $\mathrm{NPV}(r,T_h)$ is nondecreasing in $T_h$, and if $\mathrm{NPV}(r',T_h)\ge0$ with $r<r'$ then $\Psi_r(T_h)>\Psi_{r'}(T_h)\ge0$, so the set of horizons that clear the bar at $r'$ is contained in the set at $r$, and $T^*_h(r)\le T^*_h(r')$. (Checks C1 to C7.)

(b) On the admissible window $[0,\bar t]$, $\bar t:=T-t_{\rm run}-T_{\rm sea}$, the value at date 0 of a start at $t_0$ is $V(t_0)=-Ce^{-rt_0}+W(e^{-r(t_0+t_{\rm run})}-e^{-rT})/r+e^{-rT}\widetilde MW$, with $V'(t_0)=e^{-rt_0}(rC-We^{-rt_{\rm run}})$, whose sign does not depend on $t_0$. The maximum over the window is at $0$ if $We^{-rt_{\rm run}}\ge rC$ and at $\bar t$ otherwise. For $t_0\in(\bar t,T-t_{\rm run}]$ there is no terminal credit, and $V$ has the same derivative. If the yield is below $r$, $V$ increases on this interval to $V(T-t_{\rm run})=-Ce^{-r(T-t_{\rm run})}<0$, so it is negative throughout and worse than not transitioning; if the yield is at least $r$, $V$ decreases there and lies below $V(\bar t)\le V(0)$. With $h_x=0$, the value at $t_0$ is $-C+W\int_{t_{\rm run}}^\tau e^{-ru}du\le-C+W(\tau-t_{\rm run})^+$, so a nonnegative value requires $C/W\le\tau-t_{\rm run}$. With $h_x\widetilde M>0$, the value at $t_0$ is nonnegative if and only if $C/W\le[e^{-rt_{\rm run}}-e^{-r\tau}]/r+e^{-r\tau}\widetilde M$. (Checks C8 to C11.)

(c) Under the exit rule, the value of a program is the present value of its interim cash flows plus $e^{-rT}$ times the capitalized seasoned run-rate, and the second term is the same for every plan that seasons all $N$ processes by $T$. Plans therefore differ only through forgone interim savings, dual running and overhead, whose first-order approximation is \eqref{eq:ext-progcost}. A process in wave $k$ reaches run-rate at $kT_{\rm par}$ and is seasoned at $kT_{\rm par}+T_{\rm sea}$, which must not exceed $T$, so $k\le\bar n_w$. For decision speed, Proposition~\ref{prop:ext-bandwidth}(a) with $\ell_{\rm dec}=\ell_0{\rm lat}_{\mathsf o}/(1-\varkappa)$ and $q\equiv1$ gives $\hat\varkappa=1+c/2-\sqrt{c+c^2/4}$ with $c=r_d\ell_0{\rm lat}_{\mathsf o}$, and
\[
\frac{d\hat\varkappa}{dc}=\frac12\Bigl[1-\frac{1+c/2}{\sqrt{c+c^2/4}}\Bigr]<0,
\]
because $(1+c/2)^2=1+c+c^2/4>c+c^2/4$. The bound $NeT_{\rm par}/(\hat\varkappa\Phi)$ of Proposition~\ref{prop:ext-bandwidth}(b) falls as $\hat\varkappa$ rises. (Checks C12, C13.) \qed

\subsection*{Proof of Proposition~\ref{prop:reg-dip}}

(a) With $\mathrm{EB}_t=\mathrm{EB}_0+\Delta\mathrm{EB}_t>0$, $\mathrm{ND}_t\le\overline{\rm lev}\,\mathrm{EB}_t$ reads ${\rm lev}_0\mathrm{EB}_0+\Delta\mathrm{ND}_t\le\overline{\rm lev}\,\mathrm{EB}_0+\overline{\rm lev}\,\Delta\mathrm{EB}_t$, which rearranges to the stated inequality. A cost that lowers tested EBITDA by one dollar a year raises the left side by $w_{\rm cov}\overline{\rm lev}$ through $-\Delta\mathrm{EB}$, in addition to its cash in $\Delta\mathrm{ND}$; a capitalized or added-back cost enters only through $\Delta\mathrm{ND}$. (Check D1.)

(b) The public form is Proposition~\ref{prop:ext-waves}(f). For the sponsor, at wave $k$ the one-time costs of all waves so far are in net debt, $\Delta\mathrm{ND}=\mathcal T_1N_k$, and the tested EBITDA change is $w_{\rm cov}(sN_{k-1}-D(n_k))$. By (a) the covenant is $\mathcal T_1(N_{k-1}+n_k)+w_{\rm cov}\overline{\rm lev}(D(n_k)-sN_{k-1})\le\mathfrak B^{\rm cov}$, which is the stated form. In both, $\widehat D$ is convex and strictly increasing with $\widehat D(0)=0$.

(i) Wave $k$ is feasible if and only if $n_k\le\bar n(N_{k-1})$. Let $\phi(x):=x+\bar n(x)$. If $\hat s\ge0$, $\bar n$ and $\phi$ are nondecreasing. If $\hat s<0$ (only possible for the sponsor, with $\hat s\ge-\mathcal T_1$), then for $x<x'$, since $\widehat D(n)-\mathcal T_1n$ is nondecreasing, $\mathcal T_1[\bar n(x)-\bar n(x')]\le\widehat D(\bar n(x))-\widehat D(\bar n(x'))=-\hat s(x'-x)\le\mathcal T_1(x'-x)$, so $\phi(x)\le\phi(x')$. With integer waves, $\lfloor a\rfloor-\lfloor b\rfloor\le\lceil a-b\rceil$, so $x+\lfloor\bar n(x)\rfloor$ is also nondecreasing on the integers. By induction, if $N_{k-1}\le N^g_{k-1}$ then $N_k\le\min\{N,\phi(N_{k-1})\}\le\min\{N,\phi(N^g_{k-1})\}=N^g_k$. Hence $N_k<N$ for $k<\hat k$, and every feasible plan has at least $\hat k$ waves. (Check D2.)

(ii) With $\bar n(x)=(\mathfrak B+\hat sx)/\hat d$, the recursion is $N^g_k=(1+\hat s/\hat d)N^g_{k-1}+\mathfrak B/\hat d$ with $N^g_0=0$, whose solution is $(\mathfrak B/\hat s)[(1+\hat s/\hat d)^k-1]$; the waves are $(\mathfrak B/\hat d)(1+\hat s/\hat d)^{k-1}$, and $\hat k$ is the smallest $k$ at which the solution reaches $N$. $\widehat D$ is linear if and only if $D$ is, and then $\sum_kD(n_k)=d_1N$ for every plan, so \eqref{eq:ext-progcost} equals $T_{\rm par}s\sum_{k=1}^{n_w}kn_k+T_{\rm par}d_1N+n_wO_w$, with $\sum_kkn_k=\sum_{k=0}^{n_w-1}(N-N_k)$. By (i), every feasible plan has $N_k\le N^g_k$ and $n_w\ge\hat k$, so both sums are at least those of the fastest ramp. The argument holds for integer plans with the rounded-down ramp. (Checks D3 to D5.)

(iii) In the unconstrained plan of Proposition~\ref{prop:ext-waves}(b), $n_k$ is strictly decreasing while positive, so $\widehat D(n_k)$ is strictly decreasing and $\hat sN_{k-1}$ is nondecreasing ($\hat s\ge0$); the dip $\widehat D(n_k)-\hat sN_{k-1}$ is therefore strictly decreasing in $k$, and a violated budget is violated in wave 1. With $n_w$ fixed, the constrained problem minimizes a convex function subject to the convex constraints $g_k(n):=\widehat D(n_k)-\hat sN_{k-1}-\mathfrak B\le0$ and $\sum_kn_k=N$. Under Slater's condition the Kuhn--Tucker conditions hold, and for positive waves
\[
T_{\rm par}[sk+D'(n_k)]+\mu_k\widehat D'(n_k)-\hat s\sum_{j>k}\mu_j=\lambda .
\]
Subtracting the condition for $k$ from that for $k+1$ gives $T_{\rm par}[s+D'(n_{k+1})-D'(n_k)]+\mu_{k+1}\widehat D'(n_{k+1})-\mu_k\widehat D'(n_k)+\hat s\mu_{k+1}=0$. If wave $k+1$ is slack, $\mu_{k+1}=0$ and $D'(n_k)-D'(n_{k+1})=s-\mu_k\widehat D'(n_k)/T_{\rm par}$; with $D'(n)=d_1+2d_2n$ this is the stated decrement, equal to $s/(2d_2)$ if wave $k$ is slack too and negative if $\mu_k>T_{\rm par}s/\widehat D'(n_k)$. (Checks D6 to D8.)

(c) Let $f$ be the program cost, $\mathcal C:=\{n\ge0:\sum_kn_k=N\}$ and $P(n):=\sum_kg_k(n)^+$. Let $n^*$ solve the hard-budget problem with multipliers $\mu^*$; for a convex program the multiplier set does not depend on which optimum is taken, and $n^*$ minimizes $f+\sum_k\mu^*_kg_k$ over $\mathcal C$ with $\mu^*_kg_k(n^*)=0$. If $\mathfrak w T_{\rm par}\ge\max_k\mu^*_k$, then for every $n\in\mathcal C$,
\[
f(n)+\mathfrak w T_{\rm par}P(n)\ge f(n)+\sum_k\mu^*_kg_k(n)^+\ge f(n)+\sum_k\mu^*_kg_k(n)\ge f(n^*)=f(n^*)+\mathfrak w T_{\rm par}P(n^*),
\]
so $n^*$ solves the soft-budget problem. Conversely, if $n^*$ solves the soft-budget problem, the subdifferential sum rule gives $t_k\in[0,1]$, with $t_k=0$ for slack waves, such that $0\in\nabla f(n^*)+\sum_k\mathfrak w T_{\rm par}t_k\nabla g_k(n^*)+N_{\mathcal C}(n^*)$, where $N_{\mathcal C}$ is the normal cone. Then $(\mathfrak w T_{\rm par}t_k)_k$ is a multiplier vector of the hard-budget problem, and uniqueness gives $\mu^*_k=\mathfrak w T_{\rm par}t_k\le\mathfrak w T_{\rm par}$ for all $k$. So if $\mathfrak w T_{\rm par}<\max_k\mu^*_k$, $n^*$ is not a soft-budget optimum. A soft-budget optimum $\tilde n$ exists ($\mathcal C$ is compact) and satisfies $f(\tilde n)+\mathfrak w T_{\rm par}P(\tilde n)<f(n^*)$; if $P(\tilde n)=0$, $\tilde n$ would be budget-feasible with $f(\tilde n)<f(n^*)$, a contradiction, so $\tilde n$ exceeds the guided dip in some wave. For the last claim, let $p$ and $p'$ be chosen at $\mathfrak w<\mathfrak w'$ from any set of plans. Optimality gives $f(p)+\mathfrak w P(p)\le f(p')+\mathfrak w P(p')$ and $f(p')+\mathfrak w'P(p')\le f(p)+\mathfrak w'P(p)$. Adding, $(\mathfrak w'-\mathfrak w)[P(p')-P(p)]\le0$, so $P(p')\le P(p)$, and then the first inequality gives $f(p)\le f(p')$. (Checks D9 to D11.)

(d) By (i), at most $N^g_{\bar n_w}$ processes can be migrated in $\bar n_w$ waves, and a process migrated later is not seasoned by $T$, so it loses $h_xM_{\rm exit}s^{\rm E}$ of terminal value (Proposition~\ref{prop:reg-speed}(c)). Raising $\mathfrak B$ raises $\bar n(\cdot)$ and hence $N^g_{\bar n_w}$; its value is bounded by the terminal value of the processes it brings within the deadline plus the fall in program cost, which in the continuous relaxation is $\sum_k\mu^*_k$ per dollar of budget by the envelope theorem. (Checks D12 to D16.) \qed

\subsection*{Proof of Lemma~\ref{lem:reg-waterfall}}

Distributions are paid in tiers. First, investors receive their capital, up to one per dollar. Second, they receive the preferred return, up to $\Lambda_T-1$ more. Third, the general partner receives everything until it holds the share $c_{\rm cy}$ of total profit: its catch-up $G$ solves $G=c_{\rm cy}(G+\Lambda_T-1)$, so $G=c_{\rm cy}(\Lambda_T-1)/(1-c_{\rm cy})$, and the tier ends at $\mathrm{MOIC}=\Lambda_T+G=1+(\Lambda_T-1)/(1-c_{\rm cy})$. Fourth, the remainder is split $(1-c_{\rm cy},c_{\rm cy})$. Investors therefore receive $\mathrm{MOIC}$ in the first two tiers, $\Lambda_T$ in the third, and $\Lambda_T+(1-c_{\rm cy})[\mathrm{MOIC}-\Lambda_T-G]=1+(1-c_{\rm cy})(\mathrm{MOIC}-1)$ in the fourth. Fees add $c_{\rm fee}T$ to the capital they pay in per dollar invested; with fees at entry and one distribution at $T$, the net IRR solves $(1+\mathrm{IRR})^T$ equal to the net multiple. (Checks E1, E2.) \qed

\subsection*{Proof of Proposition~\ref{prop:reg-metrics}}

(a) The sponsor pays $(1+\mathfrak p)P_0$ for the equity of a firm without net debt plus costs $c_{\rm tx}P_0$, and borrows $\mathrm{ND}_0$. Net debt accrues at $r_{\rm debt}$ and is reduced by the free cash flow swept to it, $d\,\mathrm{ND}=(r_{\rm debt}\mathrm{ND}-\mathrm{FCF})\,dt$, whose solution at $T$ is the stated $\mathrm{ND}_T$. Exit equity is $M_{\rm exit}[\mathrm{EB}^{\rm sq}_T+h_x\Delta\mathrm{EB}_T]-\mathrm{ND}_T$ by the EBITDA rule, which is affine in the program's run-rate and cash flows, so the program adds $[h_xM_{\rm exit}\Delta\mathrm{EB}_T+\int_0^Te^{r_{\rm debt}(T-t)}d\,\Delta\mathrm{CF}]/\mathrm{Eq}_0$ to the multiple. With one entry and one exit, $(1+\mathrm{IRR})^T=\mathrm{MOIC}$. One percent of $\mathrm{EB}_0$ of seasoned run-rate adds $0.01h_xM_{\rm exit}\mathrm{EB}_0$ to exit equity. (Checks E3 to E5.)

(b) After-tax operating income changes by $(1-\tau)\Delta\mathrm{OI}_y$, divided by the share count; the gap to guidance follows from Assumption~\ref{ass:reg-regime}(iii). (Check E6.)

(c) Proposition~\ref{prop:tasks-aggregate}(b) is Part I's envelope result at an arbitrary weight $X$, so it holds at $X_{\mathsf o}$; only $g^*$, $\sigma^*$ and $X-k$ depend on the weight. The value is Definition~\ref{def:reg-value}. For a going concern at small share, Proposition~\ref{prop:tasks-aggregate}(c) holds at $r_{\mathsf o}$ by Proposition~\ref{prop:reg-weight}(a). (Check E7.)

(d) At given output, Lemma~\ref{lem:ext-retained} gives a price fall of $\bar\varphi$ per dollar of unit-cost saving in cost of revenue, so customers gain $\bar\varphi W_{\rm cogs}$ and pre-tax profit rises by $(1-\bar\varphi)W_{\rm cogs}+W_{\rm op}$, of which the share $\tau$ is taxed. Severance is the cash term of Assumption~\ref{ass:tasks-transition}(iii). \qed

\subsection*{Proof of Proposition~\ref{prop:reg-lbo}}

(a) The sponsor puts in $\mathrm{Eq}_0=(1+\mathfrak p)P_0+C_{\rm tx}-\mathrm{ND}_0$ and expects
\[
\mathbb E\,\mathrm{Eq}_T=M_{\rm exit}(\mathrm{EB}^{\rm sq}_T+h_x\Delta\mathrm{EB}_T)+\int_0^Te^{r_{\rm debt}(T-t)}\mathrm{FCF}\,dt-\mathrm{ND}_0e^{r_{\rm debt}T}-p_bL_b .
\]
Hence
\[
\mathrm{NPV}^{\rm sp}=e^{-r_{\rm pe}T}\mathbb E\,\mathrm{Eq}_T-\mathrm{Eq}_0=V^{\rm sp}(\Pi^{\rm pe})-\mathrm{ND}_0e^{-(r_{\rm pe}-r_{\rm debt})T}-e^{-r_{\rm pe}T}p_bL_b-(1+\mathfrak p)P_0-C_{\rm tx}+\mathrm{ND}_0,
\]
which is \eqref{eq:reg-lbo}. It is affine and decreasing in $\mathfrak p$, which gives the condition $\mathfrak p^{\rm req}\le\mathfrak p^{\max}$. Since $\mathrm{Eq}_0$ does not depend on $r_{\rm pe}$, $\mathrm{NPV}^{\rm sp}$ is strictly decreasing in $r_{\rm pe}$ when $\mathbb E\,\mathrm{Eq}_T>0$, with root $\ln(\mathbb E\,\mathrm{Eq}_T/\mathrm{Eq}_0)/T$. (Checks E3, F4.)

(b) With $r_{\rm pe}=r_{\rm debt}=r_w$, $\mathfrak L=0$ and $V^{\rm sp}=\int_0^Te^{-r_wt}\mathrm{FCF}\,dt+e^{-r_wT}\mathrm{EV}_T$. With the fair exit price $\mathrm{EV}_T=\int_T^\infty e^{-r_w(t-T)}\mathrm{FCF}(t;\Pi^{\rm pe})\,dt$, $V^{\rm sp}=\int_0^\infty e^{-r_wt}\mathrm{FCF}(t;\Pi^{\rm pe})\,dt=V^{\rm sq}+V^{\rm prog}(\Pi^{\rm pe})$, because cash flows add. Subtracting $P_0=V^{\rm sq}+V^{\rm prog}(\Pi^{\rm pub})$ gives \eqref{eq:reg-exec}. Neither the premium nor $\mathfrak L$ changes the firm's cash flows; they reallocate them among target holders, the sponsor and lenders. (Checks F1 to F3.)

(c) The public owner chooses $\Pi^{\rm pub}(\mathfrak w)$ to maximize $V^{\rm prog}(\Pi)-\mathfrak w\mathcal P(\Pi)$ over its feasible plans. The revealed-preference argument at the end of the proof of Proposition~\ref{prop:reg-dip}(c), with cost $-V^{\rm prog}$, shows that $V^{\rm prog}(\Pi^{\rm pub}(\mathfrak w))$ is nonincreasing in $\mathfrak w$, so $\Delta V^{\rm exec}$ is nondecreasing. A tighter covenant or an earlier deadline shrinks the sponsor's feasible set, so the best value over it cannot rise. For the growth weight, the sponsor picks the frontier point of $R_{X^{\rm pe}}$ (Proposition~\ref{prop:reg-weight}(c)), at which $\Delta=\rho\Delta^*(X)$ with $\rho=[(X^{\rm pe}-k)/(X-k)]^{\alpha/(1-\alpha)}$ by \eqref{eq:Dstar}, and the loss in $R_X$ follows as in the proof of Proposition~\ref{prop:reg-weight}(b). (Checks F6, F7.)

(d) Add and subtract $V^{\rm sq}+V^{\rm prog}(\Pi^{\rm pe})$ in $\mathfrak p^{\max}P_0=V^{\rm sp}(\Pi^{\rm pe})+\mathfrak L-C_{\rm tx}-e^{-r_{\rm pe}T}p_bL_b-P_0$ and use $P_0=V^{\rm sq}+V^{\rm prog}(\Pi^{\rm pub})$. (Check F5.) \qed

\subsection*{Proof of Proposition~\ref{prop:reg-hybrid}}

(a) $X^{\rm my}(\mathfrak w)=k+(X-k)/(1+\mathfrak w)$ is decreasing in $\mathfrak w$ for $X>k$. The program's cost is nondecreasing and its excess dip nonincreasing in $\mathfrak w$ by Proposition~\ref{prop:reg-dip}(c), so values at $\mathfrak w_{\rm dc}$ lie between those at $0$ and at $\mathfrak w_{\rm pub}$. (Checks G1, D10.)

(b) (i) New equity used to repay debt, or held as cash, lowers net debt by $C_P$, so by Proposition~\ref{prop:reg-dip}(a) the covenant becomes $\Delta\mathrm{ND}+\overline{\rm lev}(-\Delta\mathrm{EB})\le\mathfrak B^{\rm cov}+C_P$, and a cash budget rises by $C_P$. Operating income does not depend on how the firm is financed, so a GAAP budget is unchanged. (ii) $\omega_mR_{X^{\rm inv}}(g,m)+(1-\omega_m)R_{X^{\rm pub}}(g,m)=[\omega_mX^{\rm inv}+(1-\omega_m)X^{\rm pub}]g+m$. (Checks G2, G3.)

(c) The rebuild route costs $C^{\rm sep}$ at date 0 and $C^{\rm par,sh}$ a year for $T_{\rm par}$ and leaves the legacy cost structure, so it saves nothing. The value of automating relative to rebuilding is therefore $\mathrm{NPV}_p-[-C^{\rm sep}-C^{\rm par,sh}\mathcal A_r(T_{\rm par})]$. Without discounting, this adds $C^{\rm sep}+C^{\rm par,sh}T_{\rm par}$ to the numerator of $T^*_h$ in Proposition~\ref{prop:tasks-transition}(b) with a negative sign, which lowers $T^*_h$ by that amount divided by $\varpi_p\bar b_pW$. (Check G4.)

(d) The weight is Proposition~\ref{prop:reg-weight}(c) at $r_{\rm vc}$ with the market rule. At exit, $V_T$ of Lemma~\ref{lem:valuation2} is linear in $m$ and differentiable in $g$, with $\partial V/\partial m=S_T\mathfrak A_T$ and $\partial V/\partial g=X^{\rm scr}\,\partial V/\partial m$ by the definition of $X^{\rm scr}$, so to first order $\Delta V_T=\mathfrak A_TS_T(\Delta m_T+X^{\rm scr}\Delta g_T)$. (Checks G5, G6.) \qed
